\section{Практическая работа №1}
\label{sec:pr1}

\subsection{Выполненные задания}

В первой практической работе были реализованы функции для работы со списками, кортежами, свёртками и типом \mintinline{haskell}{Maybe}. В отдельной части были добавлены функции \mintinline{haskell}{myZipSave}, \mintinline{haskell}{myUnzipSave}, \mintinline{haskell}{myReverse}, \mintinline{haskell}{myFoldl1}, \mintinline{haskell}{myFoldr1}, \mintinline{haskell}{myTakeWhile}, \mintinline{haskell}{mySpan}, \mintinline{haskell}{myMaybe}, а также функции для пользовательского списка и примеры из лекции про приготовление тортов.

Основные коммиты практики: \mintinline{text}{f0752bd}, \mintinline{text}{5baf79c}, \mintinline{text}{78c0460}, \mintinline{text}{f476e64}, \mintinline{text}{c575c1c}, \mintinline{text}{1c028cd}, \mintinline{text}{5b5342c}.

\subsection{Использование нейросетей}

Использование нейросети было зафиксировано в комментариях к файлам практики. Модель использовалась для обсуждения структуры функций и проверки отдельных реализаций, после чего код приводился к стилю практической работы.

\subsection{Вопросы защиты и их решения}

На защите была получена задача реализовать функцию типа \mintinline{haskell}{Either (a -> a) (c -> c) -> [a] -> [c] -> [(a, c)]}. Если передан \mintinline{haskell}{Left f}, функция применяется к элементу первого списка; если \mintinline{haskell}{Right f} --- к элементу второго списка.

\captionof{listing}{Решение задачи защиты практической работы №1}
\label{lst:lab1-defense}
\begin{minted}{haskell}
foo :: Either (a -> a) (c -> c) -> [a] -> [c] -> [(a, c)]
foo _ [] _ = []
foo _ _ [] = []
foo (Left f) (x : xs) (y : ys) = (f x, y) : foo (Left f) xs ys
foo (Right f) (x : xs) (y : ys) = (x, f y) : foo (Right f) xs ys
\end{minted}

\newpage
